\scanpage{27}
$p|_{\dot X_0\times I}$ induce $\dot X_0\times I_{\unsure{2|\ell}}\longrightarrow Y_0\times I$,
being identified with $p_0\times\operatorname{id}_I$, where $p_0:X_0\longrightarrow Y_0$.

\deleted{Compatibility}\inserted{Transversality} of a \deleted{finite} family
$(B_i,E_{B_i})$ of cylindrical neighborhoods \unsure{of the|of} subsets $B_i$ of $X$.
The condition is expressed by saying that for
every $J\subset I$, $I$ finite, the structure of
\[
\bigcap_{i\in J}E_{B_i}\overset{\mathrm{def}}=E_J,
\quad\text{\unsure{around|over}}\quad
B_J\overset{\mathrm{def}}=\bigcap_{i\in J}B_i,
\]
is that of $B_J\times I^J$, the homothety
$\lambda_{E_j}$ of $E_j$ \inserted{(for $j\in J$)} inducing
\[
(b,(t_\alpha)_{\alpha\in J})\longmapsto(b,\unsure{(t'_\alpha)|(t'_\alpha]}),
\qquad
t'_\alpha=t_\alpha\ \text{if }\alpha\ne j,\quad
\text{and }\unsure{t'_\alpha|t_\alpha}=\lambda t_\alpha\ \text{if }\alpha=j.
\]
\margin{this homeomorphism $E_J\simeq B_J\times I^J$ is
canonical, if
$\dot B_J=\bigcap_{j\in J}\dot E_j
\underset{\substack{\text{\unsure{multiplied}}\\
\text{\deleted{\unsure{by}\ \illegible{1}}}\\
\text{the \unsure{unit|zero} element of}\\I^J}}{\xrightarrow{\sim}}B_J$
\ldots}

This can also be made explicit by requiring
that $E_J$ be cylindrical in each $E_{B_j}$ and so on\ldots
This transversality relation passes to germs.

One says that the preceding family \uline{and}
a conical neighborhood $T$ of $Y\subset X$ are
\inserted{transverse} \deleted{\unsure{compatible|compatible}}
if \deleted{the structure of the $T\cap E_i$}
\inserted{\uuline{in addition} $\forall\deleted{\unsure{j}}J$ as above,}
\margin{\unsure{induced by} \unsure{\(\mathinner{\cdotp-\cdotp}\)}
$E_J\simeq B_J\times I^J$ and}
\[
\begin{array}{l@{\quad}lcl@{\quad}l}
& E_J\cap T & \simeq & T_J\times I^J
& (\text{where }T_J=T\cap B_J)\\
&\cup&&\cup&\\
\text{inducing}&\dot E_J\cap\dot T&\simeq&\dot T_J\times I^J
&(\text{where }\dot T_J=\dot T\cap B_J)\\
&E_J\cap Y&\simeq&Y_J\times I^J
&(\text{where }Y_J=Y\sim B_J)
\end{array}
\]
\margin{26}
